\documentclass[12pt,a4paper]{article}

% --- Pacotes ---
\usepackage[portuges]{babel}
\usepackage[utf8]{inputenc}
\usepackage[T1]{fontenc}
\usepackage{geometry}
\geometry{margin=2.5cm}
\usepackage{longtable}
\usepackage{booktabs}
\usepackage{xcolor}
\usepackage{hyperref}
\hypersetup{
    colorlinks=true,
    linkcolor=blue!60!black,
    urlcolor=blue!60!black,
    citecolor=blue!60!black
}
\usepackage{listings}
\lstset{
    basicstyle=\ttfamily\small,
    breaklines=true,
    frame=single,
    backgroundcolor=\color{gray!5},
    rulecolor=\color{gray!30}
}
\usepackage{titlesec}
\usepackage{float}
\usepackage{caption}

\title{\textbf{Softinsa Web Server} \\ \large Documentação da API REST}
\author{Softinsa}
\date{Maio 2026}

\begin{document}

\maketitle
\thispagestyle{empty}
\newpage

\tableofcontents
\newpage

% ===================================================================
\section{Introdução}
% ===================================================================

A API do \textbf{Softinsa Web Server} é uma interface REST desenvolvida em \texttt{Node.js} com \texttt{Express 5}, utilizando \texttt{PostgreSQL} como base de dados e \texttt{Sequelize} como ORM.

\begin{itemize}
    \item \textbf{Base URL:} \texttt{http://localhost:3000/api}
    \item \textbf{Porta:} 3000
    \item \textbf{Formato de dados:} JSON
    \item \textbf{Autenticação:} JWT (Bearer token)
\end{itemize}

\subsection{Autenticação}

A API utiliza dois middleware de autenticação:

\begin{description}
    \item[\texttt{auth.middleware}] --- É aplicado globalmente a todos os pedidos. Extrai o token JWT do header \texttt{Authorization: Bearer <token>} e define \texttt{req.user}. Se não existir token válido, o utilizador é marcado como \texttt{guest}.
    \item[\texttt{requireAuth.middleware}] --- É aplicado em rotas específicas. Se \texttt{req.user.role === "guest"}, devolve \texttt{401 Unauthorized}.
\end{description}

\subsection{Papéis (Roles)}

\begin{tabular}{ll}
    \toprule
    \textbf{Código} & \textbf{Papel} \\
    \midrule
    A  & Administrador \\
    TM & Talent Manager \\
    SL & Service Line Leader \\
    CO & Consultor \\
    guest & Utilizador não autenticado \\
    \bottomrule
\end{tabular}

\newpage

% ===================================================================
\section{Rotas}
% ===================================================================

Todas as rotas são precedidas pelo prefixo \texttt{/api}.

% -------------------------------------------------------------------
\subsection{Autenticação}
\label{sec:autenticacao}
% -------------------------------------------------------------------

Prefix: \texttt{/api/autenticacao}

\begin{longtable}{@{}lllp{6cm}@{}}
    \toprule
    \textbf{Método} & \textbf{Path} & \textbf{Auth} & \textbf{Descrição} \\
    \midrule
    \endhead
    \bottomrule
    \endfoot

    POST & \texttt{/register} & --- & Regista um novo utilizador. \\
    POST & \texttt{/login} & --- & Autentica um utilizador e devolve um token JWT. \\
    GET  & \texttt{/get-me} & --- & Devolve os dados do utilizador autenticado. \\
    PUT  & \texttt{/update-me} & \texttt{requireAuth} & Edita os dados do utilizador autenticado. \\
    DELETE & \texttt{/delete-me} & \texttt{requireAuth} & Elimina o utilizador autenticado. \\
\end{longtable}

\subsubsection*{Exemplo: Register}
\begin{lstlisting}
POST /api/autenticacao/register
Body: {
    "nome": "João Silva",
    "email": "joao@exemplo.com",
    "password": "minhaPassword123"
}
Resposta (201): {
    "message": "Conta criada com sucesso",
    "user": {
        "id": 1,
        "email": "joao@exemplo.com",
        "role": "CO"
    }
}
\end{lstlisting}

\subsubsection*{Exemplo: Login}
\begin{lstlisting}
POST /api/autenticacao/login
Body: {
    "email": "joao@exemplo.com",
    "password": "minhaPassword123"
}
Resposta (200): {
    "token": "eyJhbGciOiJIUzI1NiIs...",
    "user": {
        "email": "joao@exemplo.com",
        "role": "CO"
    }
}
\end{lstlisting}

% -------------------------------------------------------------------
\subsection{Áreas}
\label{sec:areas}
% -------------------------------------------------------------------

Prefix: \texttt{/api/areas}

\begin{longtable}{@{}lllp{6cm}@{}}
    \toprule
    \textbf{Método} & \textbf{Path} & \textbf{Auth} & \textbf{Descrição} \\
    \midrule
    \endhead
    \bottomrule
    \endfoot

    GET    & \texttt{/get} & --- & Lista todas as áreas (filtra inativas para não-admin). \\
    GET    & \texttt{/:id/get} & --- & Devolve uma área pelo ID. \\
    POST   & \texttt{/create} & Admin & Cria uma nova área. \\
    PUT    & \texttt{/:id/update} & Admin & Atualiza uma área. \\
    DELETE & \texttt{/:id/delete} & Admin & Soft delete de uma área (\texttt{estado\_a\_i = false}). \\
\end{longtable}

\subsubsection*{Exemplo: Criar Área}
\begin{lstlisting}
POST /api/areas/create
Body: {
    "id_serviceline": 1,
    "nome_area": "Cloud Computing",
    "descricao_area": "Área de cloud",
    "imagem_area": "cloud.png",
    "estado_a_i": true
}
Resposta (201): { "mensagem": "Área criada" }
\end{lstlisting}

% -------------------------------------------------------------------
\subsection{Badges}
\label{sec:badges}
% -------------------------------------------------------------------

Prefix: \texttt{/api/badges}

\begin{longtable}{@{}lllp{6cm}@{}}
    \toprule
    \textbf{Método} & \textbf{Path} & \textbf{Auth} & \textbf{Descrição} \\
    \midrule
    \endhead
    \bottomrule
    \endfoot

    GET & \texttt{/get} & --- & Lista todos os badges. \\
    GET & \texttt{/:id/get} & --- & Devolve um badge pelo ID. \\
    POST & \texttt{/create} & Admin & Cria um novo badge. \\
    PUT & \texttt{/:id/update} & --- & Atualiza um badge. \\
    DELETE & \texttt{/:id/delete} & --- & Soft delete de um badge. \\
    GET & \texttt{/:id\_badge/consultor/:id\_consultor/estado} & --- & Devolve o estado de um badge para um consultor. \\
\end{longtable}

\subsubsection*{Estados possíveis do badge}
\begin{tabular}{ll}
    \toprule
    \textbf{Estado} & \textbf{Descrição} \\
    \midrule
    Concluido & Badge concluído (dentro da validade). \\
    Expirado  & Badge concluído mas expirado. \\
    Em análise & Pedido em análise (estados 1 ou 2). \\
    Por Obter & Pedido devolvido ou aguardando (estados 3, 5 ou 6). \\
    \bottomrule
\end{tabular}

\subsubsection*{Exemplo: Criar Badge}
\begin{lstlisting}
POST /api/badges/create
Body: {
    "id_area": 1,
    "nome_badge": "AWS Practitioner",
    "descricao_badge": "Badge de AWS",
    "pontos_badge": 100,
    "pago": false,
    "nivel_badge": 1,
    "sla": 30,
    "validade": 365,
    "estado_a_i": true
}
Resposta (201): { "mensagem": "Badge criado com sucesso" }
\end{lstlisting}

% -------------------------------------------------------------------
\subsection{Candidaturas}
\label{sec:candidaturas}
% -------------------------------------------------------------------

Prefix: \texttt{/api/candidaturas}

\begin{longtable}{@{}lllp{6cm}@{}}
    \toprule
    \textbf{Método} & \textbf{Path} & \textbf{Auth} & \textbf{Descrição} \\
    \midrule
    \endhead
    \bottomrule
    \endfoot

    POST & \texttt{/} & --- & Submete uma candidatura (pedido + documentos). \\
\end{longtable}

\subsubsection*{Exemplo: Submeter Candidatura}
\begin{lstlisting}
POST /api/candidaturas
Body: {
    "id_consultor": 1,
    "id_badge": 2,
    "documentos": ["certificado.pdf", "apresentacao.pdf"]
}
Resposta: { "success": true, "message": "Candidatura submetida com sucesso" }
\end{lstlisting}

O serviço de candidaturas atribui automaticamente o Talent Manager com menos pedidos e o Service Line Leader da área do badge.

% -------------------------------------------------------------------
\subsection{Consultores}
\label{sec:consultores}
% -------------------------------------------------------------------

Prefix: \texttt{/api/consultores}

\begin{longtable}{@{}lllp{6cm}@{}}
    \toprule
    \textbf{Método} & \textbf{Path} & \textbf{Auth} & \textbf{Descrição} \\
    \midrule
    \endhead
    \bottomrule
    \endfoot

    PUT & \texttt{/:id} & --- & Edita os dados de um consultor. \\
\end{longtable}

\subsubsection*{Exemplo: Editar Consultor}
\begin{lstlisting}
PUT /api/consultores/1
Body: {
    "nome": "João Atualizado",
    "email": "joao.novo@exemplo.com",
    "id_area": 2,
    "foto_perfil": "nova_foto.png",
    "password": "novaPassword"
}
Resposta: { "success": true, "message": "Dados atualizados com sucesso" }
\end{lstlisting}

% -------------------------------------------------------------------
\subsection{Learning Paths}
\label{sec:learningpaths}
% -------------------------------------------------------------------

Prefix: \texttt{/api/learningPaths}

\begin{longtable}{@{}lllp{6cm}@{}}
    \toprule
    \textbf{Método} & \textbf{Path} & \textbf{Auth} & \textbf{Descrição} \\
    \midrule
    \endhead
    \bottomrule
    \endfoot

    GET    & \texttt{/get} & --- & Lista todas as Learning Paths. \\
    GET    & \texttt{/:id/get} & --- & Devolve uma Learning Path pelo ID. \\
    POST   & \texttt{/create} & Admin & Cria uma nova Learning Path. \\
    PUT    & \texttt{/:id/update} & --- & Atualiza uma Learning Path. \\
    DELETE & \texttt{/:id/delete} & --- & Soft delete de uma Learning Path. \\
\end{longtable}

\subsubsection*{Exemplo: Criar Learning Path}
\begin{lstlisting}
POST /api/learningPaths/create
Body: {
    "nome_learning_path": "Cloud Expert",
    "descricao_learning_path": "Percurso de cloud",
    "imagem_learning_path": "cloud_path.png",
    "estado_a_i": true
}
Resposta (201): { "mensagem": "Learning Path criada com sucesso" }
\end{lstlisting}

% -------------------------------------------------------------------
\subsection{Pedidos de Badges}
\label{sec:pedidos}
% -------------------------------------------------------------------

Prefix: \texttt{/api/pedidos}

\begin{longtable}{@{}lllp{6cm}@{}}
    \toprule
    \textbf{Método} & \textbf{Path} & \textbf{Auth} & \textbf{Descrição} \\
    \midrule
    \endhead
    \bottomrule
    \endfoot

    GET  & \texttt{/get} & \texttt{requireAuth} & Lista todos os pedidos (consultor vê apenas os seus). \\
    GET  & \texttt{/:id/get} & \texttt{requireAuth} & Devolve um pedido pelo ID. \\
    POST & \texttt{/create} & \texttt{requireAuth} & Cria um novo pedido (apenas CO ou Admin). \\
    POST & \texttt{/:id/tm-review} & --- & Avaliação do Talent Manager (aprovar/devolver). \\
    POST & \texttt{/:id/sl-review} & --- & Avaliação do Service Line Leader (aprovar/devolver/rejeitar). \\
    POST & \texttt{/:id/resubmit} & --- & Reenvio do pedido pelo consultor. \\
\end{longtable}

\subsubsection*{Workflow de Pedidos}
O fluxo de aprovação segue os seguintes estados:

\begin{tabular}{ll}
    \toprule
    \textbf{Estado} & \textbf{Descrição} \\
    \midrule
    1 & Pedido criado (aguarda TM) \\
    2 & Reenviado pelo consultor \\
    3 & Devolvido pelo TM \\
    4 & --- \\
    5 & Aprovado pelo TM \\
    6 & Devolvido pelo SL \\
    7 & Rejeitado pelo SL \\
    8 & Aprovado / Concluído \\
    \bottomrule
\end{tabular}

\subsubsection*{Fluxo}
\begin{lstlisting}[frame=none]
  Consultor  -->  TM  -->  SL  -->  Concluído (8)
      |            |        |
      |        devolver   rejeitar (7)
      |         (3)       devolver (6)
      |            |        |
      +<-----------+--------+
      resubmit (2)
\end{lstlisting}

\subsubsection*{Exemplo: Criar Pedido}
\begin{lstlisting}
POST /api/pedidos/create
Body: { "id_consultor": 1, "id_badge": 2 }
Resposta (201): {
    "mensagem": "Pedido criado com sucesso",
    "dados": { "id_pedido_badge": 1, ... }
}
\end{lstlisting}

\subsubsection*{Exemplo: TM Review}
\begin{lstlisting}
POST /api/pedidos/1/tm-review
Body: { "acao": "aprovar" }
Resposta (200): { "mensagem": "Pedido aprovado" }
\end{lstlisting}

\subsubsection*{Exemplo: SL Review}
\begin{lstlisting}
POST /api/pedidos/1/sl-review
Body: { "acao": "aprovar" }
Resposta (200): { "mensagem": "Pedido aprovado com sucesso" }
\end{lstlisting}

\subsubsection*{Exemplo: Reenviar Pedido}
\begin{lstlisting}
POST /api/pedidos/1/resubmit
Resposta (200): { "mensagem": "Pedido reenviado com sucesso" }
\end{lstlisting}

% -------------------------------------------------------------------
\subsection{Service Lines}
\label{sec:servicelines}
% -------------------------------------------------------------------

Prefix: \texttt{/api/serviceLines}

\begin{longtable}{@{}lllp{6cm}@{}}
    \toprule
    \textbf{Método} & \textbf{Path} & \textbf{Auth} & \textbf{Descrição} \\
    \midrule
    \endhead
    \bottomrule
    \endfoot

    GET    & \texttt{/get} & --- & Lista todas as Service Lines. \\
    GET    & \texttt{/:id/get} & --- & Devolve uma Service Line pelo ID. \\
    POST   & \texttt{/create} & Admin & Cria uma nova Service Line. \\
    PUT    & \texttt{/:id/update} & --- & Atualiza uma Service Line. \\
    DELETE & \texttt{/:id/delete} & --- & Soft delete de uma Service Line. \\
\end{longtable}

\subsubsection*{Exemplo: Criar Service Line}
\begin{lstlisting}
POST /api/serviceLines/create
Body: {
    "id_learning_path": 1,
    "nome_serviceline": "Cloud Services",
    "descricao_serviceline": "Linha de cloud",
    "imagem_serviceline": "cloud_sl.png",
    "estado_a_i": true
}
Resposta (201): { "mensagem": "Service Line criada com sucesso" }
\end{lstlisting}

% -------------------------------------------------------------------
\subsection{Utilizadores}
\label{sec:utilizadores}
% -------------------------------------------------------------------

Prefix: \texttt{/api/utilizadores}

\begin{longtable}{@{}lllp{6cm}@{}}
    \toprule
    \textbf{Método} & \textbf{Path} & \textbf{Auth} & \textbf{Descrição} \\
    \midrule
    \endhead
    \bottomrule
    \endfoot

    GET    & \texttt{/get} & \texttt{requireAuth} & Lista todos os utilizadores (consultor vê apenas o seu perfil). \\
    GET    & \texttt{/:id/get} & --- & Devolve um utilizador pelo ID. \\
    POST   & \texttt{/create} & \texttt{requireAuth} & Cria um novo utilizador com papel específico. \\
    PUT    & \texttt{/:id/update} & --- & Atualiza um utilizador. \\
    DELETE & \texttt{/:id/delete} & Admin & Elimina um utilizador. \\
    POST   & \texttt{/:id/objetivo/create} & --- & Adiciona um objetivo a um utilizador. \\
    DELETE & \texttt{/:id/objetivo/delete} & --- & Remove um objetivo de um utilizador. \\
    GET    & \texttt{/:idUtilizador/notificacoes} & \texttt{requireAuth} & Lista notificações de um utilizador. \\
    POST   & \texttt{/:idUtilizador/notificacoes/create} & \texttt{requireAuth} & Envia uma notificação a um utilizador. \\
\end{longtable}

\subsubsection*{Tipos de Utilizador}
\begin{tabular}{ll}
    \toprule
    \textbf{Código} & \textbf{Papel} \\
    \midrule
    c  & Consultor \\
    tm & Talent Manager \\
    sl & Service Line Leader \\
    a  & Admin \\
    \bottomrule
\end{tabular}

\subsubsection*{Exemplo: Criar Utilizador}
\begin{lstlisting}
POST /api/utilizadores/create
Body: {
    "nome_utilizador": "Maria Santos",
    "email_utilizador": "maria@exemplo.com",
    "password_utilizador": "hash123",
    "username_utilizador": "maria.santos",
    "tipo_utilizador": "c",
    "imagem_utilizador": "avatar.png",
    "id_area": 1
}
Resposta (201): {
    "success": true,
    "message": "Utilizador criado com sucesso",
    "id_utilizador": 5
}
\end{lstlisting}

\newpage

% ===================================================================
\section{Modelos de Dados}
% ===================================================================

A base de dados é composta pelas seguintes tabelas (modelos Sequelize):

\begin{longtable}{@{}ll@{}}
    \toprule
    \textbf{Tabela} & \textbf{Descrição} \\
    \midrule
    \endhead
    \bottomrule
    \endfoot

    Utilizadores & Utilizadores do sistema (todos os papéis). \\
    Administradores & Administradores. \\
    Consultores & Consultores (com área e pontuação). \\
    TalentManagers & Talent Managers. \\
    ServiceLineLiders & Service Line Leaders. \\
    Areas & Áreas de conhecimento. \\
    ServiceLines & Linhas de serviço. \\
    LearningPaths & Percursos de aprendizagem. \\
    Badges & Badges disponíveis. \\
    BadgesConcluidos & Badges concluídos por consultores. \\
    Objetivos & Objetivos definidos para consultores. \\
    PedidosBadges & Pedidos de badges. \\
    HistoricoPedidos & Histórico de estados dos pedidos. \\
    Estados & Tabela de estados. \\
    NotificacoesPedidos & Notificações relacionadas com pedidos. \\
    NotificacoesAdmin & Notificações de administrador. \\
    Documentacoes & Documentos associados a pedidos. \\
    Enviadas & Registos de envios. \\
    Favoritos & Favoritos dos utilizadores. \\
    Conquistas & Conquistas do sistema. \\
    ConquistasConsultores & Conquistas atribuídas a consultores. \\
    Politicas & Políticas/RGPD. \\
    PoliticasAceites & Políticas aceites por utilizadores. \\
    Requisitos & Requisitos de badges. \\
    EstadoPedidos & Estados de pedidos. \\
    RGPD & Registos RGPD. \\
\end{longtable}

\newpage

% ===================================================================
\section{Rotas Não Implementadas / Descontinuadas}
% ===================================================================

As seguintes rotas existem nos ficheiros de código mas estão comentadas ou por implementar:

\begin{itemize}
    \item \textbf{Gestão} (\texttt{Gestao.route.js}) --- Rotas para SLA e RGPD. Não ativa no router principal.
    \item \textbf{Conquistas} (\texttt{Conquistas.route.js}) --- CRUD de conquistas. Comentada (tabela Conquistas pode não existir).
    \item Histórico de pedido: \texttt{GET /api/pedidos/:id/historico} --- por implementar.
    \item Atualizar pedido: \texttt{PUT /api/pedidos/:id/update} --- por implementar.
    \item Eliminar pedido: \texttt{DELETE /api/pedidos/:id/delete} --- por implementar.
    \item Requisitos de badges --- todo o subgrupo de rotas comentado.
    \item Dashboard: \texttt{GET /api/utilizadores/:id/dashboard} --- por implementar.
    \item Badges do utilizador: \texttt{GET /api/utilizadores/:id/badges} --- por implementar.
\end{itemize}

\end{document}
